\documentclass[10pt,a4paper]{article}
\usepackage[utf8]{inputenc}
\usepackage[margin=0.60in]{geometry}
\usepackage{hyperref}
\usepackage{enumitem}
\usepackage{titlesec}
\usepackage{xcolor}
\definecolor{linkblue}{RGB}{0,0,139}
\definecolor{highlightred}{RGB}{155,43,43}
\hypersetup{colorlinks=true,linkcolor=linkblue,urlcolor=linkblue,filecolor=linkblue,pdfauthor={Siyuan Mei},pdftitle={Curriculum Vitae - Siyuan Mei}}
\titleformat{\section}{\large\bfseries}{}{0em}{}[\titlerule]
\titlespacing{\section}{0pt}{9pt}{4pt}
\setlist[itemize]{leftmargin=1.25em,itemsep=1pt,topsep=1pt,parsep=0pt}
\setlist[enumerate]{leftmargin=1.7em,itemsep=2pt,topsep=1pt,parsep=0pt}
\setlength{\parindent}{0pt}\setlength{\parskip}{2pt}\setlength{\emergencystretch}{2em}
\newcommand{\entry}[3]{\textbf{#1}\hfill #2\\[-1pt]#3}
\newcommand{\pub}[4]{\item[\textbf{#1}] #2. \textit{#3}. #4.}
\newcommand{\tagline}[1]{\textcolor{highlightred}{\textbf{#1}}}
\begin{document}
\begin{center}
    {\fontsize{30}{34}\selectfont\bfseries Siyuan Mei}\\[6pt]
    \href{mailto:siyuan.mei@fau.de}{siyuan.mei@fau.de}
    \enspace | \enspace \href{https://siyuan-mei.github.io}{Homepage (github.io)}
    \enspace | \enspace \href{https://scholar.google.com/citations?user=W_xPF2YAAAAJ\&hl=en}{Google Scholar}\\[3pt]
    \href{https://lme.tf.fau.de/persons/siyuan-mei/}{Lab page}
    \enspace | \enspace \href{https://www.linkedin.com/in/siyuan-mei-409353247/}{LinkedIn}
    \enspace | \enspace \href{https://github.com/siyuan-mei}{GitHub}\\[4pt]
    Pattern Recognition Lab (LME), FAU Erlangen-Nürnberg, Germany
\end{center}

\section{RESEARCH INTERESTS}
\textbf{Medical AI:} reconstruction, synthesis, segmentation, and registration.\\
\textbf{Computer vision:} diffusion and flow-matching models, generative modeling, representation learning, and dense prediction.\\
\textbf{LLMs \& software engineering:} LLM-assisted development, reproducible research tooling, and agentic systems.

\section{EDUCATION}
\entry{FAU Erlangen-Nürnberg}{2023--2027 (expected)}{Ph.D. in Computer Science; Supervisor: Prof. Andreas Maier.}\\[4pt]
\entry{Peking University Health Science Center}{May--Oct. 2026}{Visiting Scholar, Artificial Intelligence for Medical Imaging (AIMI) Lab; Supervisor: Prof. Yixing Huang.}\\[4pt]
\entry{FAU Erlangen-Nürnberg}{2020--2023}{M.Sc. in Medical Engineering; thesis grade: 1.0.}\\[4pt]
\entry{Zhejiang University of Science and Technology}{2016--2020}{B.Eng. in Automation; Sino-German 2+3 Program; honored graduate scholarship.}

\section{PROFESSIONAL EXPERIENCE}
\textbf{Friedrich-Alexander-Universität Erlangen-Nürnberg (FAU)} \hfill 2023--2027\\
\textit{Ph.D. Student, Pattern Recognition Lab; Supervisor: Prof. Andreas Maier}
\begin{itemize}
    \item Fully funded through the ERC project \textit{Advancing osteoporosis medicine by observing bone microstructure and remodelling using a 4D nanoscope}; developed \textbf{BigReg}, an efficient registration pipeline for high-resolution X-ray microscopy and light-sheet fluorescence microscopy.
    \item Conduct research in medical AI and computer vision, adapting foundation models and developing medical image generation methods with GANs, diffusion models, and flow matching.
    \item Teach \textit{Flat-panel CT Reconstruction} and contribute lecture slides for \textit{Vibe Coding: Software Engineering with LLMs}, a Springer book project.
    \item \textbf{Outcomes:} first-author work in Pattern Recognition and MICCAI; a Featured Paper in Journal of Medical Imaging; and an Oral Presentation at MICCAI 2026.
\end{itemize}
\textbf{Peking University Health Science Center, AIMI Lab} \hfill May--Oct. 2026\\
\textit{Visiting Scholar; Supervisor: Prof. Yixing Huang}\\
Artificial Intelligence for Medical Imaging Lab, Institute of Medical Technology, 240 Pathology Building, No. 38 Xueyuan Road, Haidian District, Beijing, China.
\begin{itemize}
    \item Participated in the MICCAI BIC-MAC Challenge and trained a multimodal-to-CT generative model for PET attenuation correction; achieved first place on the validation-set CT metric.
    \item Built practical knowledge of clinical CT applications and imaging workflows; visited the ICT Center in Chongqing and attended MISC 2026 at Tsinghua University.
\end{itemize}

\section[SELECTED PUBLICATIONS]{SELECTED PUBLICATIONS \hfill \textnormal{\textcolor{linkblue}{J=Journal, C=Conference, M=Manuscript}}}
\begin{enumerate}
    \pub{[C.1]}{\textbf{Siyuan Mei} et al}{Time Matters: Rethinking Diffusion and Flow Models in One-Step Medical Image Translation}{MICCAI 2026; \tagline{Early Accept}; \tagline{Oral Presentation}; first author}
    \pub{[C.2]}{Mengchen Fan et al., including Siyuan Mei}{Detail Consistent Stage-Wise Distillation for Efficient 3D MRI Segmentation}{MICCAI 2026; \tagline{Early Accept}; \tagline{Spotlight Presentation}}
    \pub{[C.3]}{Daiqi Liu et al., including Siyuan Mei}{Speech-Guided Multimodal Learning for Vocal Tract Segmentation in Real-Time MRI}{MICCAI 2026; \tagline{Oral Presentation}}
    \pub{[C.4]}{\textbf{Siyuan Mei} et al}{WING: A Window-Prior-Based Generative Network with Gated Inception for Cross-Modality CT Synthesis}{MICCAI MIART Workshop 2026; accepted workshop paper; first author}
    \pub{[J.1]}{\textbf{Siyuan Mei} et al}{Vision Transformer Hook for Dense Predictions}{Pattern Recognition 179 (2026), 113818; first author}
    \pub{[J.2]}{\textbf{Siyuan Mei} et al}{BigReg: An Efficient Registration Pipeline for High-Resolution X-Ray and Light-Sheet Fluorescence Microscopy}{Journal of Medical Imaging 12(5), 054004 (2025); \tagline{Featured Paper}}
    \pub{[C.5]}{\textbf{Siyuan Mei}, F. Fan, A. Maier}{Metal Inpainting in CBCT Projections Using Score-based Generative Model}{IEEE ISBI 2023; first author}
    \pub{[C.6]}{\textbf{Siyuan Mei} (co-author)}{TCCT: Trajectory-Conditioned CBCT Reconstruction for Sinusoidal Non-Circular Orbits with a Fourier Neural Operator}{MICCAI 2026; Poster}
    \pub{[J.3]}{Y. Sun et al., including Siyuan Mei}{Filter2Noise: A Framework for Interpretable and Zero-Shot Low-Dose CT Image Denoising}{Journal of Medical Imaging 13(2), 024004 (2026)}
    \pub{[J.4]}{C. Ye et al., including Siyuan Mei}{Robustness and Stability Analysis of Differentiable Shift-Variant FBP for Cone-Beam CT under Challenging Acquisition Settings}{Machine Learning for Biomedical Imaging (2026)}
    \pub{[J.5]}{M. Thies et al., including Siyuan Mei}{A Gradient-Based Approach to Fast and Accurate Head Motion Compensation in Cone-Beam CT}{IEEE Transactions on Medical Imaging (2024)}
    \pub{[J.6]}{C. Ye et al., including Siyuan Mei}{DRACO: Differentiable Reconstruction for Arbitrary CBCT Orbits}{Physics in Medicine and Biology 70, 075005 (2025)}
    \pub{[C.7]}{M. Gu et al., including Siyuan Mei}{Unsupervised Domain Adaptation Using Soft-Labeled Contrastive Learning with Reversed Monte Carlo Method for Cardiac Image Segmentation}{MICCAI 2024}
    \pub{[C.8]}{L. Pfaff et al., including Siyuan Mei}{No-New-Denoiser: A Critical Analysis of Diffusion Models for Medical Image Denoising}{MICCAI 2024}
    \pub{[M.1]}{Zirong Li, \textbf{Siyuan Mei} et al}{Generative Drifting for Conditional Medical Image Generation}{arXiv preprint (2026)}
    \pub{[M.2]}{\textbf{Siyuan Mei}, Y. Xia, F. Fan}{GANeXt: A Fully ConvNeXt-Enhanced Generative Adversarial Network for MRI- and CBCT-to-CT Synthesis}{arXiv preprint (2025)}
\end{enumerate}

\section{HONORS AND AWARDS}
\begin{itemize}
    \item \textbf{1st place, Task 2} and \textbf{4th place, Task 1}, MICCAI SynthRAD2025 Challenge, 2025.
    \item \textbf{Two merit-based scholarships} during M.Sc. studies at FAU Erlangen-Nürnberg.
    \item Fully funded Ph.D. through the ERC 4D Nanoscope project, FAU, 2023--2027.
    \item Honored graduate scholarship, Zhejiang University of Science and Technology, 2020.
\end{itemize}

\section{TEACHING}
\entry{Flat-panel CT Reconstruction}{FAU}{Summer semesters 2024--2026.}\\[3pt]
\entry{Vibe Coding: Software Engineering with LLMs}{FAU}{Summer semester 2026; lecture slides accompanying a Springer book.}\\[4pt]
\section{ACADEMIC SERVICES}
\textbf{Journal reviewer:} Pattern Recognition; Journal of Machine Learning for Biomedical Imaging (MELBA).\\
\textbf{Technical editor:} Journal of Medical Imaging.\\
\textbf{Conference reviewer:} MICCAI.

\section{ADDITIONAL INFORMATION}
\textbf{Languages:} Chinese (native), English (professional), German (C1 certificate).\\
\textbf{Interests:} fitness and tennis (level 2.5). \quad \textbf{MBTI (self-reported):} ENTJ-A.\\
\textbf{Career conversations:} I am open to research conversations, collaborations, and job-related discussions in medical imaging, computer vision, and LLM-assisted software engineering; I appreciate hearing about suitable opportunities.\\
\textbf{Open-source and teaching resources:} \href{https://github.com/siyuan-mei/WING}{WING}, \href{https://github.com/siyuan-mei/JiR}{JiR}, \href{https://github.com/siyuan-mei/ViT-Hook}{ViT-Hook}, and \href{https://github.com/PRLab-FAU/VibeCodingSlides}{Vibe Coding lecture slides} accompanying a Springer book.
\end{document}
